\appendix
\chapter{Python Implementation}
\label{appendix:code}

This appendix reproduces the complete Python implementation of the
Monte Carlo simulation engine described in Sections~8--10. The code
is structured in seven modules, each corresponding to a distinct
layer of the framework. All results reported in the paper are fully
reproducible by running \texttt{main()} with the default
\texttt{ModelParams} instance and \texttt{seed = 42}.

The implementation requires Python 3.10 or later and the following
standard scientific libraries: \texttt{numpy}, \texttt{scipy}, and
\texttt{matplotlib}. No proprietary or non-standard packages are used.

\section*{Module 0 — Dependencies and Style}

Global imports and \texttt{matplotlib} configuration. The colour
palette uses Barrick Gold's corporate navy and gold tones for
consistency with the figures in the main text.

\begin{lstlisting}[style=python, caption={Dependencies and global style configuration}]
import numpy as np
import matplotlib.pyplot as plt
import matplotlib.patches as mpatches
from scipy import stats
from dataclasses import dataclass
from typing import Tuple
import warnings
warnings.filterwarnings("ignore")

plt.rcParams.update({
    "font.family": "serif",
    "axes.spines.top": False,
    "axes.spines.right": False,
    "axes.grid": True,
    "grid.alpha": 0.3,
    "figure.dpi": 150,
})

COLORS = {
    "primary":   "#1B3A6B",   # deep navy
    "secondary": "#C9A84C",   # gold (Barrick reference)
    "accent":    "#8B1A1A",   # deep red
    "light":     "#D6E4F0",
    "mid":       "#7FB3D3",
}
\end{lstlisting}

\section*{Module 1 — Parameter Dataclass}

All model inputs are collected in a single \texttt{ModelParams}
dataclass. This design makes the framework portable: switching from
the Base scenario to the Bear or Bull scenario (Section~10) requires
only overriding the relevant fields, leaving the simulation
architecture unchanged. Default values correspond to the Base
scenario calibrated in Section~8.6.

\begin{lstlisting}[style=python, caption={ModelParams dataclass — all calibrated inputs}]
@dataclass
class ModelParams:
    """
    All model parameters in one place.
    Sources: Barrick Gold Annual Report 2023, Damodaran (2024),
             LBMA gold price series 2000-2024.
    """
    # Simulation
    N_sim:      int   = 20_000    # number of Monte Carlo draws
    seed:       int   = 42        # reproducibility

    # Projection horizon
    T_high:     int   = 5         # years in high-growth phase
    T_trans:    int   = 5         # years in transition phase

    # Gold price GBM
    P0:         float = 2_050.0   # USD/oz  (spot, end-2023)
    mu_P:       float = 0.03      # annual drift
    sigma_P:    float = 0.155     # annual volatility (hist. 2000-2024)

    # Production volumes
    Q0:         float = 4.05      # Moz  (Barrick 2023 production)
    Q_mature:   float = 3.80      # Moz  (conservative long-run)
    sigma_Q:    float = 0.06      # annual vol around guidance

    # Operating economics
    aisc_start: float = 0.62      # AISC / Revenue  (2023 actual ~62%)
    aisc_mature:float = 0.65      # slight cost pressure at maturity
    da_pct:     float = 0.08      # D&A as % of revenue
    capex_pct_high:  float = 0.12 # net capex / revenue  high-growth
    capex_pct_stable:float = 0.09 # net capex / revenue  stable
    tax_rate:   float = 0.30      # effective corporate tax rate

    # ROIC and growth (used for TV identity)
    ROIC_high:  float = 0.12
    ROIC_stable:float = 0.10
    g_high:     float = 0.05
    g_trans:    float = 0.025
    g_stable:   float = 0.025     # terminal growth rate (<= risk-free)

    # WACC -- Ornstein-Uhlenbeck process
    r0:         float = 0.086     # starting WACC
    r_bar:      float = 0.090     # long-run equilibrium WACC
    kappa:      float = 0.40      # mean-reversion speed
    sigma_r:    float = 0.012     # WACC volatility

    # Correlation structure
    rho:        float = -0.25     # corr(dW^P, dW^r)

    # Balance sheet (equity bridge)
    net_debt:   float = 4_900.0   # USD million
    shares:     float = 1_690.0   # million shares outstanding
\end{lstlisting}

\section*{Module 2 — Stochastic Path Generation}

\texttt{StochasticPaths} generates $N$ joint realisations of the
three stochastic drivers. The price and WACC Brownian motions are
correlated via a Cholesky decomposition of the $2\times 2$
correlation matrix (Section~8.5). Both the GBM and the
Ornstein-Uhlenbeck discretisations are \emph{exact}: no Euler
approximation error is introduced at any step size $\Delta t$.

\begin{lstlisting}[style=python, caption={StochasticPaths — exact GBM and OU discretisations}]
class StochasticPaths:
    """
    Generates N_sim joint realisations of:
      - Gold price path {P_t}  via exact GBM discretisation  (eq. 79)
      - WACC path       {r_t}  via exact OU discretisation   (eq. 80)
      - Volume shocks   {Q_t}  via log-normal perturbation   (eq. 70)
    All arrays have shape (N_sim, T).
    """

    def __init__(self, p: ModelParams):
        self.p = p
        self.T = p.T_high + p.T_trans
        rng = np.random.default_rng(p.seed)

        # Cholesky decomposition for correlated Brownians
        corr_matrix = np.array([[1.0, p.rho],
                                 [p.rho, 1.0]])
        L   = np.linalg.cholesky(corr_matrix)
        raw = rng.standard_normal((p.N_sim, self.T, 2))
        Z   = raw @ L.T                          # shape (N, T, 2)
        Z_Q = rng.standard_normal((p.N_sim, self.T))

        self.P = self._gbm_paths(Z[:, :, 0], p)
        self.r = self._ou_paths(Z[:, :, 1], p)
        self.Q = self._volume_paths(Z_Q, p)

    @staticmethod
    def _gbm_paths(Z: np.ndarray, p: ModelParams) -> np.ndarray:
        """Exact GBM discretisation -- eq. (79)."""
        dt = 1.0
        log_increments = (p.mu_P - 0.5 * p.sigma_P**2) * dt \
                         + p.sigma_P * np.sqrt(dt) * Z
        log_P = np.cumsum(log_increments, axis=1)
        return p.P0 * np.exp(log_P)

    @staticmethod
    def _ou_paths(Z: np.ndarray, p: ModelParams) -> np.ndarray:
        """Exact OU discretisation -- eq. (80)."""
        dt      = 1.0
        T       = p.T_high + p.T_trans
        e_kdt   = np.exp(-p.kappa * dt)
        std_cond = p.sigma_r * np.sqrt(
            (1 - np.exp(-2 * p.kappa * dt)) / (2 * p.kappa))
        r = np.empty((p.N_sim, T + 1))
        r[:, 0] = p.r0
        for t in range(T):
            r[:, t+1] = (p.r_bar
                         + (r[:, t] - p.r_bar) * e_kdt
                         + std_cond * Z[:, t])
        return r[:, 1:]

    @staticmethod
    def _volume_paths(Z_Q: np.ndarray, p: ModelParams) -> np.ndarray:
        """
        Log-normal volume shock around deterministic guidance -- eq. (70).
        E[Q_t] = Q_bar_t by construction.
        """
        T     = p.T_high + p.T_trans
        t_idx = np.arange(1, T + 1)
        Q_bar = p.Q0 + (p.Q_mature - p.Q0) * (t_idx / T)
        shocks = np.exp(-0.5 * p.sigma_Q**2 + p.sigma_Q * Z_Q)
        return Q_bar[np.newaxis, :] * shocks
\end{lstlisting}

\section*{Module 3 — DCF Engine}

\texttt{DCFEngine} computes FCFF and enterprise value for all $N$
simulation draws simultaneously using vectorised NumPy operations —
no explicit Python loop over simulations. The terminal value uses
the ROIC--growth identity of equation~(48); the denominator is
floored at 50 basis points to prevent numerical instability in draws
where $r_T \approx g_s$.

\begin{lstlisting}[style=python, caption={DCFEngine — vectorised FCFF computation and terminal value}]
class DCFEngine:
    """
    Computes FCFF and Enterprise Value for each simulation draw.
    Fully vectorised over all N_sim paths.
    """

    def __init__(self, paths: StochasticPaths, p: ModelParams):
        self.paths = paths
        self.p     = p
        self.T     = p.T_high + p.T_trans
        self.N     = p.N_sim

    def _revenue(self) -> np.ndarray:
        return self.paths.P * self.paths.Q             # (N, T)

    def _aisc_ratio(self) -> np.ndarray:
        t = np.arange(1, self.T + 1)
        return (self.p.aisc_start
                + (self.p.aisc_mature - self.p.aisc_start) * (t / self.T))

    def _capex_ratio(self) -> np.ndarray:
        t = np.arange(1, self.T + 1)
        return (self.p.capex_pct_high
                + (self.p.capex_pct_stable - self.p.capex_pct_high)
                * (t / self.T))

    def compute_fcff(self) -> np.ndarray:
        """FCFF = NOPAT - net capex  (eq. 73)."""
        Rev   = self._revenue()
        EBIT  = Rev * (1 - self._aisc_ratio() - self.p.da_pct)
        NOPAT = EBIT * (1 - self.p.tax_rate)
        capex = Rev * self._capex_ratio()
        return NOPAT - capex                           # (N, T)

    def compute_discount_factors(self) -> np.ndarray:
        """DF_t = 1 / prod_{s=1}^{t}(1+r_s)  (eq. 75)."""
        return 1.0 / np.cumprod(1 + self.paths.r, axis=1)

    def compute_terminal_value(self, df: np.ndarray) -> np.ndarray:
        """TV using ROIC-growth identity (eq. 48), discounted to t=0."""
        p        = self.p
        Rev_T    = self._revenue()[:, -1]
        EBIT_T   = Rev_T * (1 - p.aisc_mature - p.da_pct)
        NOPAT_T1 = EBIT_T * (1 - p.tax_rate) * (1 + p.g_stable)
        fcff_T1  = NOPAT_T1 * (1 - p.g_stable / p.ROIC_stable)
        r_T      = self.paths.r[:, -1]
        denom    = np.maximum(r_T - p.g_stable, 0.005)  # floor at 50 bps
        return (fcff_T1 / denom) * df[:, -1]

    def compute_ev(self) -> Tuple[np.ndarray, np.ndarray, np.ndarray]:
        """Returns ev, pv_fcff, pv_tv -- all shape (N,)."""
        fcff    = self.compute_fcff()
        df      = self.compute_discount_factors()
        pv_fcff = np.sum(fcff * df, axis=1)
        pv_tv   = self.compute_terminal_value(df)
        ev      = pv_fcff + pv_tv
        return ev, pv_fcff, pv_tv
\end{lstlisting}

\section*{Module 4 — Monte Carlo Simulator}

\texttt{MonteCarloSimulator} orchestrates path generation, DCF
computation, and the equity bridge from enterprise value to implied
share price. The \texttt{summary()} method returns the full set of
distributional statistics reported in Tables~9--13.

\begin{lstlisting}[style=python, caption={MonteCarloSimulator — orchestration and equity bridge}]
class MonteCarloSimulator:
    """Orchestrates path generation, DCF computation, output collection."""

    def __init__(self, params: ModelParams = None):
        self.p = params or ModelParams()
        paths  = StochasticPaths(self.p)
        engine = DCFEngine(paths, self.p)
        self.ev, self.pv_fcff, self.pv_tv = engine.compute_ev()
        self.paths = paths
        # equity bridge
        self.equity_value    = self.ev - self.p.net_debt
        self.price_per_share = self.equity_value / self.p.shares

    def summary(self) -> dict:
        ev = self.ev
        ps = self.price_per_share
        return {
            "EV_mean":    np.mean(ev),
            "EV_median":  np.median(ev),
            "EV_std":     np.std(ev),
            "EV_p5":      np.percentile(ev,  5),
            "EV_p25":     np.percentile(ev, 25),
            "EV_p75":     np.percentile(ev, 75),
            "EV_p95":     np.percentile(ev, 95),
            "SE":         np.std(ev) / np.sqrt(self.p.N_sim),
            "TV_share":   np.mean(self.pv_tv / ev),
            "Price_mean": np.mean(ps),
            "Price_p5":   np.percentile(ps,  5),
            "Price_p95":  np.percentile(ps, 95),
        }
\end{lstlisting}

\section*{Module 5 — Sensitivity Analysis}

Spearman rank correlations between each stochastic driver and the
resulting enterprise value, as reported in Table~10. Rank
correlations are preferred over Pearson correlations because they are
robust to the non-linear, non-Gaussian relationship between gold
prices and EV.

\begin{lstlisting}[style=python, caption={Sensitivity analysis via Spearman rank correlations}]
def sensitivity_analysis(sim: MonteCarloSimulator) -> dict:
    """
    Spearman rank correlation between each driver and EV.
    Higher |rho| -> greater influence on enterprise value.
    """
    ev       = sim.ev
    drivers  = {
        "Gold Price (terminal)": sim.paths.P[:, -1],
        "Gold Price (average)":  sim.paths.P.mean(axis=1),
        "Production Volume":     sim.paths.Q[:, -1],
        "WACC (terminal)":       sim.paths.r[:, -1],
        "WACC (average)":        sim.paths.r.mean(axis=1),
    }
    results = {}
    for name, driver in drivers.items():
        rho, pval = stats.spearmanr(driver, ev)
        results[name] = {"rho": rho, "pval": pval}
    return results
\end{lstlisting}

\section*{Module 6 — Visualisation}

Four plot functions produce the figures embedded in the paper.
Each function accepts an optional \texttt{ax} argument for embedding
in multi-panel layouts.

\begin{lstlisting}[style=python, caption={Visualisation functions — EV distribution, fan charts, tornado}]
def plot_ev_distribution(sim: MonteCarloSimulator, ax=None):
    """Histogram + KDE of the EV distribution with percentile bands."""
    s  = sim.summary()
    ev = sim.ev / 1_000
    if ax is None:
        fig, ax = plt.subplots(figsize=(8, 4.5))
    ax.hist(ev, bins=120, density=True, color=COLORS["light"],
            edgecolor="white", linewidth=0.3)
    kde     = stats.gaussian_kde(ev)
    x_range = np.linspace(ev.min(), ev.max(), 500)
    ax.plot(x_range, kde(x_range), color=COLORS["primary"], linewidth=2)
    p5, p25, p75, p95 = (np.percentile(ev, q) for q in [5, 25, 75, 95])
    ax.axvspan(p5,  p95,  alpha=0.10, color=COLORS["primary"], label="P5-P95")
    ax.axvspan(p25, p75,  alpha=0.18, color=COLORS["primary"], label="P25-P75")
    ax.axvline(s["EV_mean"]   / 1_000, color=COLORS["primary"],
               linestyle="--", linewidth=1.5,
               label=f'Mean: ${s["EV_mean"]/1000:.1f}B')
    ax.axvline(s["EV_median"] / 1_000, color=COLORS["secondary"],
               linestyle="-",  linewidth=1.5,
               label=f'Median: ${s["EV_median"]/1000:.1f}B')
    ax.set_xlabel("Enterprise Value (USD billion)")
    ax.set_ylabel("Density")
    ax.legend(fontsize=9)
    return ax


def plot_price_fan_chart(sim: MonteCarloSimulator, ax=None):
    """Fan chart of gold price paths -- GBM percentile bands."""
    if ax is None:
        fig, ax = plt.subplots(figsize=(8, 4.5))
    T     = sim.p.T_high + sim.p.T_trans
    years = np.arange(0, T + 1)
    P_all = np.hstack([np.full((sim.p.N_sim, 1), sim.p.P0), sim.paths.P])
    P_pct = np.percentile(P_all, [5, 25, 50, 75, 95], axis=0)
    ax.fill_between(years, P_pct[0], P_pct[4], alpha=0.15,
                    color=COLORS["secondary"], label="P5-P95")
    ax.fill_between(years, P_pct[1], P_pct[3], alpha=0.30,
                    color=COLORS["secondary"], label="P25-P75")
    ax.plot(years, P_pct[2], color=COLORS["secondary"],
            linewidth=2, label="Median")
    ax.set_xlabel("Year"); ax.set_ylabel("Gold Price (USD/oz)")
    ax.legend(fontsize=9); ax.set_xticks(years)
    return ax


def plot_wacc_fan_chart(sim: MonteCarloSimulator, ax=None):
    """Fan chart of WACC paths -- OU percentile bands."""
    if ax is None:
        fig, ax = plt.subplots(figsize=(8, 4.5))
    T     = sim.p.T_high + sim.p.T_trans
    years = np.arange(1, T + 1)
    r_pct = np.percentile(sim.paths.r * 100, [5, 25, 50, 75, 95], axis=0)
    ax.fill_between(years, r_pct[0], r_pct[4], alpha=0.15,
                    color=COLORS["primary"], label="P5-P95")
    ax.fill_between(years, r_pct[1], r_pct[3], alpha=0.30,
                    color=COLORS["primary"], label="P25-P75")
    ax.plot(years, r_pct[2], color=COLORS["primary"],
            linewidth=2, label="Median WACC")
    ax.axhline(sim.p.r_bar * 100, color="grey", linewidth=0.8,
               linestyle=":",
               label=f"Equilibrium ({sim.p.r_bar*100:.1f}%)")
    ax.set_xlabel("Year"); ax.set_ylabel("WACC (%)")
    ax.legend(fontsize=9); ax.set_xticks(years)
    return ax


def plot_tornado(sens: dict, ax=None):
    """Tornado chart of Spearman rank correlations."""
    if ax is None:
        fig, ax = plt.subplots(figsize=(7, 4))
    names = list(sens.keys())
    rhos  = [sens[n]["rho"] for n in names]
    order = np.argsort(np.abs(rhos))
    ns    = [names[i] for i in order]
    rs    = [rhos[i]  for i in order]
    cols  = [COLORS["primary"] if r > 0 else COLORS["accent"] for r in rs]
    ax.barh(ns, rs, color=cols, edgecolor="white", height=0.55)
    ax.axvline(0, color="black", linewidth=0.8)
    ax.set_xlabel("Spearman Rank Correlation with EV")
    pos = mpatches.Patch(color=COLORS["primary"], label="Positive impact")
    neg = mpatches.Patch(color=COLORS["accent"],  label="Negative impact")
    ax.legend(handles=[pos, neg], fontsize=9)
    return ax
\end{lstlisting}

\section*{Module 7 — Main Entry Point}

The \texttt{main()} function runs the full pipeline: simulation,
summary statistics, sensitivity analysis, and figure generation. All
output files are written to the working directory.

\begin{lstlisting}[style=python, caption={Main entry point — full pipeline execution}]
def main():
    params = ModelParams()
    sim    = MonteCarloSimulator(params)
    s      = sim.summary()
    sens   = sensitivity_analysis(sim)

    # Print summary
    print(f"Mean EV:    ${s['EV_mean']:,.0f}M")
    print(f"Median EV:  ${s['EV_median']:,.0f}M")
    print(f"P5 / P95:   ${s['EV_p5']:,.0f}M  /  ${s['EV_p95']:,.0f}M")
    print(f"TV share:   {s['TV_share']*100:.1f}%")
    print(f"Std Error:  ${s['SE']:,.0f}M")

    # Four-panel figure
    fig, axes = plt.subplots(2, 2, figsize=(14, 10))
    fig.suptitle("Three-Stage Stochastic DCF -- Barrick Gold",
                 fontsize=14, fontweight="bold")
    plot_ev_distribution(sim, ax=axes[0, 0])
    plot_price_fan_chart(sim, ax=axes[0, 1])
    plot_wacc_fan_chart(sim,  ax=axes[1, 0])
    plot_tornado(sens,        ax=axes[1, 1])
    plt.tight_layout()
    plt.savefig("mc_dcf_results.png", dpi=180, bbox_inches="tight")
    plt.close()

    return sim, s, sens


if __name__ == "__main__":
    sim, s, sens = main()
\end{lstlisting}
